% ==================== 第 1 章 选题背景与研究意义 ====================
\clearpage
\section{选题背景与研究意义}

\subsection{问题背景}

个性化辅导排课并不是简单地把课程填入日历，而是一个多资源、多约束、
动态变化的组合优化问题。每个课次至少涉及课程、班级、教师、教室和时间槽
五类对象；对象之间又存在资质、容量、设备、不可用时间和偏好时段等关系。
当课次数量增加或教师临时请假时，局部调整可能触发连锁变更，人工逐项核对
容易遗漏冲突，也难以稳定复现同一套决策规则。

现有教务流程常以电子表格维护基础数据，再通过即时通信工具确认变更。该
方式适合小规模、低频变化场景，但存在三方面局限：一是约束分散在人员经验
和多个表格中，难以形成统一规则；二是方案质量依赖人工搜索顺序，无法保证
硬约束始终满足；三是方案调整缺少结构化差异记录，解释、审批和通知成本较高。

因此，本课题选择“约束求解 + 独立校验 + 可解释交互”的技术路线，把不可
违反的业务规则与可权衡的运营偏好分开建模，以机器求解负责组合搜索，以
独立校验器复核结果，以可视化界面呈现方案差异。

\subsection{研究意义}

\subsubsection{方法意义}

课题将教师、教室、班级和时间槽冲突统一表示为离散约束，并把偏好命中、
资源利用、负载均衡和变更数量纳入同一多目标框架。与只输出单一课表相比，
这种建模方式能够明确区分“必须满足”和“尽量满足”，便于分析不同权重对
方案的影响，也便于在业务规则变化时扩展模型。

\subsubsection{应用意义}

系统可将重复核对工作转化为数据校验和自动求解，帮助教务人员把精力集中到
规则确认、异常处理和方案审批。独立校验器能够对输出课表再次检查，降低因
程序实现或接口传输造成的冲突风险；局部重排机制则能在教师请假等情况下
优先保留原方案，减少通知范围。

\subsubsection{平台意义}

飞书多维表格适合维护结构化业务对象，消息卡片和审批流程适合承接方案确认
与通知。本课题先完成可独立运行的排课核心，再设计飞书集成接口，可避免把
关键求解逻辑绑定在单一前端，同时为后续形成“数据维护—智能排课—审批发布—
变更通知”的协同闭环奠定基础。

\subsection{应用场景与研究边界}

本课题聚焦 K12 一对一及小班辅导的周排课与临时调课。研究对象包括教师、
教室、班级、课程请求和离散时间槽；输出为满足全部硬约束的课表及其评分、
解释和变更明细。第一阶段不讨论学费结算、招生 CRM、在线授课、考勤和复杂
权限体系，也不把自然语言模型作为排课正确性的来源。自然语言只作为未来的
输入辅助，最终约束仍需转化为可校验的结构化数据。
